\section{Discrete scale invariance and \texorpdfstring{$1/f$}{1/f} spectra}
\label{app:one_over_f}
The ``pink'' ($1/f$) noise ansatz used in Part~VI can be derived from a minimal scale-invariance mechanism: a log-uniform distribution of characteristic times (continuous scale invariance) or a geometric ladder of times (discrete scale invariance). For background on $1/f$ noise phenomenology and standard mechanisms (including broad mixtures of relaxation times), see, e.g., \cite{DuttaHorn1981,Weissman1988}.

\begin{proposition}[$1/f$ spectrum from a log-uniform mixture of Lorentzians]
\label{prop:one_over_f}
Let $x(t)$ be a stationary process obtained as a superposition of independent Ornstein--Uhlenbeck components with relaxation times $\tau\in[\tau_{\min},\tau_{\max}]$ whose density is log-uniform,
\begin{equation}
    \rho(\tau)=\frac{C}{\tau}\,\mathbf{1}_{[\tau_{\min},\tau_{\max}]}(\tau),
\end{equation}
and whose single-component power spectral density is Lorentzian, $S_\tau(f)=\frac{A\,\tau}{1+(2\pi f\tau)^2}$. Then the total PSD satisfies
\begin{equation}
    S(f)=\int_{\tau_{\min}}^{\tau_{\max}} \rho(\tau)\,S_\tau(f)\,d\tau \;=\; \frac{A C}{2\pi}\,\frac{1}{f}\int_{2\pi f\tau_{\min}}^{2\pi f\tau_{\max}} \frac{du}{1+u^2}.
\end{equation}
In the intermediate band $1/\tau_{\max}\ll f\ll 1/\tau_{\min}$, the integral is approximately constant, hence $S(f)\propto 1/f$.
\end{proposition}
\begin{proof}[Reference]
This follows by the change of variables $u=2\pi f\tau$ and estimating the arctangent integral in the intermediate band.
\end{proof}

\paragraph{Discrete scale invariance.}
If the characteristic times form a geometric ladder $\tau_k=\tau_0\,\phi^k$ (golden scaling) with weights approximating a log-uniform density, the resulting spectrum exhibits the same $1/f$ envelope with log-periodic modulations (a periodic function of $\log_\phi f$), consistent with the ``log-periodic peaks'' template used in Part~VI \cite{Sornette1998DSI}.

\begin{proposition}[Common-mode field implies positive cross-correlation]
\label{prop:common_mode_crosscorr}
Let two colocated detectors have outputs
\begin{equation}
    y_1(t)=x(t)+n_1(t),\qquad y_2(t)=x(t)+n_2(t),
\end{equation}
where $x(t)$ is a stationary ``common'' signal and $n_1,n_2$ are stationary detector noises that are uncorrelated with $x$ and with each other. Then the cross power spectral density satisfies
\begin{equation}
    S_{12}(f)\equiv \langle \widetilde{y}_1(f)\,\overline{\widetilde{y}_2(f)}\rangle = S_x(f),
\end{equation}
so the cross-spectrum inherits the sign/shape of the common field and is strictly positive for a positive PSD $S_x(f)$.
\end{proposition}
\begin{proof}[Reference]
This is the standard cross-spectrum identity for a shared common-mode signal with uncorrelated noises.
\end{proof}

